\chapter{El Metadata Store}
\label{ch:metadata}

Los metadatos por arista (pares clave-valor arbitrarios en
\code{HashMap<String, String>} que los parsers llenan) son el
campo de tamaño más variable asociado a una relación. Inlinearlos
dentro de \Compact{} forzaría un registro de largo variable y
destruiría el invariante de 32 bytes
(Capítulo~\ref{ch:compact-relation}). El motor en cambio guarda
los metadatos externamente en \code{MetadataStore} y deja sólo un
offset de 64 bits en \Compact{}.

\section{Estructura}

\code{MetadataStore} es un buffer de bytes append-only con una
codificación simple prefijada por longitud:

\begin{lstlisting}[style=rust]
pub struct MetadataStore {
    buf: Vec<u8>,   // append-only
}
// layout de un record en offset o:
//   u32 num_entries
//   num_entries veces:
//     u16 key_len, bytes de la key, u16 value_len, bytes del value
\end{lstlisting}

El primer offset de record válido es 1, así que
$\code{metadata\_offset} = 0$ en \Compact{} significa
inequívocamente ``sin metadatos''. En el \code{append}, el store
devuelve el offset donde escribió el record nuevo y ese offset
vive en \Compact{} el resto de la sesión.

\section{Operaciones}

\begin{algorithm}
\caption{\textsc{Metadata-Append}}\label{alg:meta-append}
\KwIn{store $M$, mapa $K : \{(k, v)\}$}
\KwOut{offset $o \geq 1$, o 0 si $K$ está vacío}
\If{$K$ está vacío}{\Return 0\;}
$o \gets M.\text{buf.len()} + 1$\;
escribir \code{u32(len($K$))} y luego cada par \code{(key, value)}
con prefijos de largo \code{u16} en $M.\text{buf}$\;
\Return $o$\;
\end{algorithm}

\begin{algorithm}
\caption{\textsc{Metadata-Read}}\label{alg:meta-read}
\KwIn{store $M$, offset $o$}
\KwOut{\code{HashMap<String, String>}, o vacío si $o = 0$}
\If{$o = 0$}{\Return mapa vacío\;}
posicionar cursor en $o - 1$ dentro de $M.\text{buf}$\;
leer \code{u32 num}\;
\For{$i \gets 0$ \KwTo $num-1$}{
    leer \code{u16 k\_len}, siguientes \code{k\_len} bytes como
    UTF-8 $\to k$\;
    leer \code{u16 v\_len}, siguientes \code{v\_len} bytes como
    UTF-8 $\to v$\;
    insertar $(k, v)$\;
}
\Return mapa\;
\end{algorithm}

\begin{complexity}
\textsc{Metadata-Append}: $\Ord{|K| + \sum |k| + \sum |v|}$.
\textsc{Metadata-Read}: lo mismo. El espacio por arista es cero
si el parser no produjo metadatos, si no la longitud en bytes del
record codificado. Una arista Sysmon típica produce $\sim$60
bytes de metadatos, comparado con los 32 bytes de \Compact{} ---
por eso importa el store externo: las aristas se mantienen
densas en el hot path del matcher y los metadatos sólo se pagan
cuando un analista los pide.
\end{complexity}

\section{Inmutabilidad}

\begin{invariant}[Metadata append-only]
Una vez que un record fue escrito en \code{MetadataStore.buf},
ninguna operación trunca o reescribe los bytes previos. Los
offsets son entonces estables por la vida del grafo. En
particular, el matcher puede capturar un offset dentro de un
\Compact{} en tiempo de ingesta y leer los metadatos a través de
él en tiempo de hunt sin preocuparse por mutaciones concurrentes
del writer.
\end{invariant}

Cuando se persiste una sesión el \code{MetadataStore.buf} entero
se serializa junto con las listas de entidades/relaciones; al
cargar, el store se reconstituye verbatim preservando todos los
offsets.

\begin{inthecode}
\filepath{graph\_hunter\_core/src/metadata\_store.rs} --- struct
y rutinas de append/read.\\
\filepath{graph\_hunter\_core/src/graph.rs:246} --- donde
\textsc{Add-Relation} llama a \code{meta\_store.append}.
\end{inthecode}
